\chapter{绪论}
\section{引言（一级标题）}
人……

基于……
\subsection{学习的定义（二级标题）}
由于
\subsection{三种学习方式}
\section{基于核的学习}
\subsection{核方法的实质}
可以从如下几个方面来理解核代入方法的实质：

首先，由T. M. Cover的模式可分性定理可以知道，一个复杂的模式分类问题映射到高维空间后，会比在低维空间里更容易变得线性可分。（A complex pattern classification problem cast in a high-dimensional space nonlinearly is more likely to be linearly separable than in a low-dimensional space.）\textsuperscript{[1]}。图1.4给出了一个低维空间中的非线性问题映射到高维空间后变成线性问题的一个示例。
\subsection{典型的基于核的学习算法}
\subsection{核学习算法的应用}
\section{本文的主要研究工作}
\section{本文的内容安排}
………

………
